\documentclass[aspectratio=169]{beamer}

\usetheme{colorful-dream}

\usepackage{minted}
\setminted{fontsize=\footnotesize, breaklines=true}

\title{Podstawy Pythona dla Data Science}
\subtitle{Wykład 1: Podstawy Pythona, środowisko, pierwszy wykres inżynierski}
\author{Lu Pa}
\date{\today}

% --- Title Slide Logos ---
\titlegraphic{
  \vspace{0.5cm}
  \begin{center}
    \includegraphics[height=1.8cm]{logo_university_placeholder.png}
    \hspace{1cm}
    \includegraphics[height=1.8cm]{logo_faculty_placeholder.png}
  \end{center}
}

\begin{document}

% ---------------------------------------------------------
\begin{frame}
  \titlepage
\end{frame}

% ---------------------------------------------------------
\section{PART I}
\begin{frame}[plain]
  \begin{center}
    \vspace{2cm}
    {\usebeamerfont{title}\color{white}\Huge PART I}
  \end{center}
\end{frame}

% ---------------------------------------------------------
\begin{frame}{Krótka historia Pythona}
  \begin{itemize}
    \item Stworzony przez Guido van Rossuma w 1991 roku
    \item Zaprojektowany jako język czytelny, prosty i ekspresyjny
    \item Inspirowany językami ABC, Modula-3 oraz paradygmatem funkcyjnym
    \item Python 2 → 2000, Python 3 → 2008 (obecny standard)
    \item Dziś: jeden z najpopularniejszych języków na świecie
  \end{itemize}
\end{frame}

% ---------------------------------------------------------
\begin{frame}{Python w liczbach}
  \begin{itemize}
    \item \textbf{\#1} najpopularniejszy język w wielu rankingach branżowych
    \item Ponad \textbf{400 000} pakietów w repozytorium PyPI
    \item Używany przez:
      \begin{itemize}
        \item Google, Meta, Microsoft, Tesla, NASA, CERN
        \item branże: automotive, lotnictwo, finanse, robotyka, biotechnologia
      \end{itemize}
    \item Dominujący w:
      \begin{itemize}
        \item data science
        \item machine learning
        \item automatyzacji i testowaniu
        \item obliczeniach naukowych
      \end{itemize}
  \end{itemize}
\end{frame}

% ---------------------------------------------------------
\begin{frame}{Dlaczego Python? (Zalety)}
  \begin{itemize}
    \item Bardzo czytelna składnia — niski próg wejścia
    \item Ogromny ekosystem bibliotek naukowych
    \item Świetne wsparcie społeczności
    \item Działa na Windows, Linux, macOS
    \item Integruje się z C/C++ dla wysokiej wydajności
    \item Idealny do prototypowania i produkcji
  \end{itemize}
\end{frame}

% ---------------------------------------------------------
\begin{frame}{Ograniczenia Pythona}
  \begin{itemize}
    \item Nie jest najszybszym językiem (interpretowany)
    \item GIL ogranicza wielowątkowość obliczeń CPU-bound
    \item Zarządzanie pakietami bywa trudne dla początkujących
    \item Nie nadaje się do systemów bardzo niskopoziomowych
    \item Jednak: biblioteki naukowe rozwiązują większość problemów wydajności
  \end{itemize}
\end{frame}

% ---------------------------------------------------------
\begin{frame}{Python vs MATLAB / C++ / R}
  \begin{itemize}
    \item \textbf{MATLAB}: świetny dla inżynierów, ale drogi i zamknięty
    \item \textbf{C++}: bardzo szybki, ale złożony i mało wygodny
    \item \textbf{R}: mocny w statystyce, słabszy w ogólnej inżynierii
    \item \textbf{Python}: darmowy, elastyczny, ogromny ekosystem, standard branżowy
  \end{itemize}
\end{frame}

% ---------------------------------------------------------
\begin{frame}{Ekosystem naukowy Pythona}
  \begin{itemize}
    \item \textbf{NumPy} — szybkie tablice numeryczne
    \item \textbf{Pandas} — tabele danych, czyszczenie, grupowanie
    \item \textbf{Matplotlib} — wykresy
    \item \textbf{SciPy} — przetwarzanie sygnałów, optymalizacja
    \item \textbf{scikit-learn} — uczenie maszynowe
    \item \textbf{Jupyter} — interaktywne notatniki
  \end{itemize}
\end{frame}

% ---------------------------------------------------------
\begin{frame}{Dlaczego Jupyter w inżynierii?}
  \begin{itemize}
    \item Kod + opis + wykresy w jednym miejscu
    \item Idealny do eksperymentów i dokumentacji
    \item Powtarzalne analizy
    \item Świetny do nauczania i raportowania
    \item Powszechnie używany w przemyśle i badaniach
  \end{itemize}
\end{frame}

% ---------------------------------------------------------
\section{PART II}
\begin{frame}[plain]
  \begin{center}
    \vspace{2cm}
    {\usebeamerfont{title}\color{white}\Huge PART II}
  \end{center}
\end{frame}

% ---------------------------------------------------------
\begin{frame}[fragile]{Pierwsze importy}
\begin{minted}{python}
import numpy as np
import pandas as pd
import matplotlib.pyplot as plt
\end{minted}
\end{frame}

% ---------------------------------------------------------
\begin{frame}{Zbiór danych}
  \begin{itemize}
    \item Syntetyczne dane telemetryczne pojazdu
    \item Kolumny:
      \begin{itemize}
        \item time\_s — czas w sekundach
        \item speed\_kmh — prędkość pojazdu
        \item accel\_mps2 — przyspieszenie
        \item engine\_rpm — obroty silnika
        \item fuel\_rate\_gps — zużycie paliwa
      \end{itemize}
  \end{itemize}
\end{frame}

% ---------------------------------------------------------
\begin{frame}[fragile]{Wczytywanie danych}
\begin{minted}{python}
df = pd.read_csv("../datasets/vehicle_speeds_trackA.csv")
df.head()
\end{minted}
\end{frame}

% ---------------------------------------------------------
\begin{frame}[fragile]{Pierwszy wykres inżynierski}
\begin{minted}{python}
plt.figure(figsize=(10,4))
plt.plot(df["time_s"], df["speed_kmh"])
plt.xlabel("Czas [s]")
plt.ylabel("Prędkość [km/h]")
plt.grid(True)
plt.title("Prędkość pojazdu w czasie (Track A)")
plt.show()
\end{minted}
\end{frame}

% ---------------------------------------------------------
\begin{frame}{Zadania praktyczne}
  \begin{itemize}
    \item Przelicz prędkość z km/h na m/s
    \item Oblicz min, max i średnią prędkość
    \item Wykreśl RPM w funkcji czasu
    \item Zaznacz na wykresie maksymalną prędkość
    \item Odfiltruj dane dla prędkości > 50 km/h
    \item Wczytaj Track B i porównaj profile jazdy
  \end{itemize}
\end{frame}

% ---------------------------------------------------------
\end{document}
